\section*{De la clausura genealógica a la realización física}

La tesis de esta obra adquiere cuatro clausuras coordinadas y no reductibles:
genealógica, operatorial, matricial y física. No se trata de yuxtaponer una
teoría de conjuntos, una teoría de álgebras de operadores o un formalismo
cuántico a una construcción numérica previa. Los cuatro desarrollos nacen de
la misma precedencia causal: APP construye la doble realización posicional y
aritmética; el TRIT orientado distingue retorno, presente torsional y apertura;
el TPK conserva valor, hoja, orientación, residuo, cociente, acarreo, borde,
ruta, registro, memoria, supervivencia, fase y calibre. Sólo después se forman
la sección compatible, sus lectores y sus representaciones. El diagrama rector
es, por ello,
\[
\begin{gathered}
 \mathrm{APP}
 \longrightarrow \mathrm{TRIT}_{\rm orientado}
 \longrightarrow \mathrm{TPK}
 \longrightarrow \mathcal X_\infty^{\rm cal},\\[2mm]
 \mathcal X_\infty^{\rm cal}
 \xrightarrow{\;\mathfrak R\;}\mathbb R,
 \qquad
 \mathcal X_\infty^{\rm cal}\times\mathcal C_{\rm exc}
 \xrightarrow{\;\mathfrak E\;}\mathcal E_{\rm exc},\\[2mm]
 C(X_m)\xrightarrow{\;\operatorname{diag}\;}M_{9^m}(\mathbb C)
 \xrightarrow{\;a\mapsto a\otimes I_9\;}
 \mathrm{UHF}(9^\infty)
 \xrightarrow{\;\pi_\tau(\,\cdot\,)''\;}R_{\rm hyp},\\[2mm]
 \mathsf P_{\rm CT}^{+}
 \xrightarrow{\;\mathfrak R_{\rm MD}^{\rm disc,+}\;}
 \mathbf{SpecRoute}\times\mathsf{Hilb}_{\mathbb R,J}
 \times\mathbf B\mathfrak M_D^+\times\mathbf B G_C.
\end{gathered}
\]
Las cuatro líneas no forman una única cadena lineal. La evaluación real y la
proyección excepcional son lectoras hermanas de la sección enriquecida; la
rama tracial pasa por las álgebras cilíndricas y sus inclusiones unitarias; la
realización física parte de rutas preparadas provistas de las estructuras que
exige cada componente. Esta separación de dominios impide deducir el factor de
von Neumann directamente de una sección, identificar la proyección
excepcional con una acción sobre cifras o presentar una magnitud física como
si fuese un estado TPK sin morfismo de realización.

\subsection*{Clausura genealógica: extensionalidad, infinito y decisión}

La primera clausura determina qué significa conservar un objeto a través de
todas las profundidades. Dado un sistema inverso de espacios finitos
\((X_m,p_m)\), una sección \(x=(x_m)_m\) no queda caracterizada únicamente por
la igualdad de sus valores publicados. El teorema de descenso por fibras fija
las condiciones bajo las cuales un lector finito desciende a un cociente y
separa sobreyectividad, constancia en fibras y unicidad del mapa descendido.
El lema de K\"onig garantiza existencia de una rama cuando los niveles son
finitos, no vacíos y finamente ramificados; la unicidad exige, además, el
carácter completo que selecciona una única prolongación compatible.

Sobre una sección fijada, los cilindros decrecientes proporcionan una
realización exacta del paso \(0\!-1\!-\!\infty\): el soporte se contrae, la
altura aumenta, la masa permanece unitaria y las medidas cilíndricas convergen
débil-* a \(\delta_x\). Este resultado coordina el delta de Kronecker finito,
las esperanzas condicionales entre niveles y el delta de Dirac como medida
puntual, sin identificar este último con un operador diferencial. La
fundación de prefijos impide cadenas descendentes infinitas en cada
presentación finita; la coinducción controla la prolongación compatible. Los
ordinales de von Neumann se elevan con una fibra de memoria de ruta, y la
cardinalidad aparece como decategorización que no retiene toda la genealogía.

La misma distinción obliga a separar potencia cilíndrica, potencia cerrada y
potencia plena, así como elección conjuntista, elección natural, elección
algebraica y elección efectiva. La codificación en ZF/ZFC conserva las
construcciones y permite contrastarlas con extensionalidad, Fundación,
Potencia y Elección; no convierte a la recta real ni al universo conjuntista
en origen de APP--TRIT--TPK\@. En el dominio efectivo, la decidibilidad de cada
ventana finita, la existencia del límite, la unicidad de una instancia y la
imposibilidad de un decisor total uniforme son cuatro afirmaciones distintas.
La reducción de G\"odel--Turing afecta sólo a la cuarta: no invalida selectores
particulares ya construidos. Análogamente, el contraste G\"odel--Cohen separa
la extensión \(M\mapsto M[G]\) del enriquecimiento interno de un estado y
evita confundir el cruce de densos de profundidad con genericidad sobre un
modelo. Los cilindros Borel, las jerarquías de comparación y el paso
proyectivo de minimax se introducen con sus hipótesis de compacidad,
continuidad y consistencia, no como consecuencias automáticas de la notación
inversa.

Esta clausura posee también una formulación informacional. Si una actualización
\(U\) del estado completo es biyectiva, entonces
\(H(X\mid U(X))=0\), y el término lógico ideal de Landauer se anula. Si un
lector \(q\) conserva sólo la sombra visible, la información descartada se
localiza en \(H(X\mid q(X))\). En la formulación operatorial, un automorfismo
traza-preservante conserva la entropía, mientras una esperanza condicional
traza-preservante \(E_N:M\to N\) satisface
\[
 S_\tau(E_Nh)-S_\tau(h)=D_\tau(h\Vert E_Nh)\geq0.
\]
El coste corresponde al olvido de la fibra, no al transporte reversible del
estado enriquecido; esta identidad no afirma disipación física total nula ni
complejidad computacional reducida.

\subsection*{Clausura operatorial: de la unidad nonádica a la dimensión continua}

La segunda clausura realiza operatorialmente la conservación evolutiva de la
unidad. Para el sistema completo se consideran
\[
 A_m=M_{9^m}(\mathbb C),
 \qquad \iota_m(a)=a\otimes I_9,
 \qquad
 \tau_m(a)=9^{-m}\operatorname{Tr}(a).
\]
La identidad \(\tau_{m+1}\circ\iota_m=\tau_m\) conserva simultáneamente
unidad y traza. El límite inductivo es la UHF de tipo sobrenatural
\(9^\infty=3^\infty\), y la clausura débil de su representación GNS tracial es
el factor hiperinfinito \(R_{\rm hyp}\) de tipo \(II_1\). Los valores
\(\operatorname{rank}(P)/9^m\) se densifican y la traza de proyectores en
\(R_{\rm hyp}\) realiza todo el intervalo \([0,1]\): una dimensión continua
surge de una torre de inclusiones discretas que preservan la unidad.

Esta construcción no se confunde con la rama marcada
\(j_{m,r}(a)=a\otimes p_r\). Las isometrías de una fibra individual generan en
norma los compactos sobre el límite hilbertiano y, en clausura débil, un factor
de tipo \(I_\infty\). La selección de un sector conserva la marca, pero no la
unidad de la torre; la elevación simultánea de las nueve fibras es la que
produce la rama UHF--\(II_1\). Tampoco se identifica sin prueba la realización
UHF con la realización por producto cruzado del reloj profinito: una
conjugación canónica requiere declarar acción, medida, amenabilidad,
ergodicidad e isotropía. La lógica ortomodular de proyectores, los contextos
booleanos, los estados normales y la normalización entrópica
\[
 S_\tau(9^m\rho)=S_{\rm vN}(\rho)-m\log9
\]
se derivan dentro de esta construcción operatorial y mantienen separados el
estado matricial, su densidad tracial y su publicación estadística.

Topología, dimensión y composición no se colapsan en el factor. La clausura
transversal se registra mediante
\[
 \mathfrak C_{\rm HMT}
 =\bigl(C(\Omega_\infty),R_{\rm hyp},\mathcal G_{\rm TPK}\bigr).
\]
El espectro conmutativo conserva puntos y cilindros; el factor tracial conserva
dimensión y estadística; el grupoide conserva concatenación, inversión
admisible y procedencia de rutas. Esta arquitectura triple se formula después
de la doble proyección y del corredor excepcional. No los funda ni los
reemplaza: la evaluación arquimediana y la cadena
Paley--Witt--Golay--Leech--\(V^\natural\)--Monstruo--Moonshine permanecen como
las dos lecturas coordinadas del mismo antecedente enriquecido.

\subsection*{Clausura matricial: profundidad, tamaño exterior y dilatación}

La tercera clausura introduce un control simultáneo de tres índices. Si
\(m\) mide profundidad genealógica, \(n\) el tamaño exterior de una cuadrícula
y \(s\) el nivel matricial de sus entradas, la familia
\(\mathfrak M_{m,s}^{(n)}\) es bigraduada por \((m,s)\) cuando \(n\) permanece
fijo y trigraduada por \((m,n,s)\) cuando también varía el tamaño exterior.
En particular,
\[
 \mathcal K_{m,s}
 :=\operatorname{UCP}\!\bigl(C(X_m),M_s(\mathbb C)\bigr)
\]
posee dos operaciones independientes: la restricción genealógica
\(\rho_m^s(\Phi)=\Phi\circ\jmath_m\) y la compresión matricial
\[
 \operatorname{mc}_{\boldsymbol V}(\Phi_1,\ldots,\Phi_k)(f)
 =\sum_rV_r^*\Phi_r(f)V_r,
 \qquad \sum_rV_r^*V_r=I_s.
\]
Ambas conmutan exactamente:
\[
 \rho_m^s\!\left(\operatorname{mc}_{\boldsymbol V}(\Phi_1,\ldots,\Phi_k)\right)
 =\operatorname{mc}_{\boldsymbol V}\!\left(
   \rho_m^{s_1}(\Phi_1),\ldots,\rho_m^{s_k}(\Phi_k)\right).
\]
La identidad es functorial y distingue refinamiento de historia de reducción
del observador.

Toda familia compatible en profundidad determina una única aplicación UCP
sobre \(C(X_\infty)\) y admite una dilatación de Stinespring común. En el
corredor diagonal, una única medida espectral proyectiva \(P\) satisface
\[
 E_x^{(m)}=V^*P([x]_m)V
\]
para todas las profundidades y todos los cilindros. La dilatación mínima es
única hasta equivalencia unitaria y conserva exactamente lo visto por el
lector; una suma directa permite una realización ambiental fiel que conserva
todas las distinciones del álgebra. Sobre el límite UHF, una familia compatible
de aplicaciones UCP posee asimismo una prolongación y una dilatación comunes.
La esperanza condicional
\[
 \mathbb E_m=\operatorname{id}\otimes\tau_9:
 M_{9^{m+1}}(\mathbb C)\longrightarrow M_{9^m}(\mathbb C)
\]
preserva la traza y descarta el último factor de forma controlada; su ambiente
de Stinespring registra el dígito omitido. En general no existe una memoria
auxiliar finita uniforme para todas las profundidades. Además, la fidelidad de
una representación no implica por sí sola fidelidad dinámica APP--TPK: ésta
exige un entrelazador compatible con ruta, monodromía y composición.

La teoría matricial se concreta en objetos verificables. Los cuadrados mágicos
cuánticos qutrit de parámetros \((9,3)\) y \((12,3)\) poseen respectivamente
81 y 144 efectos, sumas de filas y columnas exactamente unitarias y
dilataciones simultáneas declaradas; las mismas POVM producen cubos mágicos
cuánticos de órdenes nueve y doce. El caso de orden doce se organiza mediante
la fibración \(12\to4\times3\), con su calibre y su condición explícita de
equivarianza. Por otra vía, las esperanzas condicionales suma y producto de APP
generan el álgebra
\[
 C^*(P_{\Sigma,2},P_{\Pi,2})
 \cong\mathbb C^4\oplus M_2(\mathbb C)\oplus M_2(\mathbb C),
\]
y el bloque central realiza en cada nivel \(s\) el intervalo matricial libre
\([5I_s,9I_s]\). Estas construcciones distinguen cuadrados latinos cuánticos,
cuadrados mágicos, medidas proyectivas, medidas positivas y semiclasicidad;
ninguna identificación se obtiene sólo porque los objetos compartan orden
nueve.

El antecedente espectral común de dos lectores se formula únicamente después
de haber construido ambos. Para mapas finitos
\(R_m:X_m\to\mathcal A_m\) y \(Q_m:X_m\to\mathcal E_m\), las incidencias
\[
 P_{a,e}=1_{\{R_m=a\}}1_{\{Q_m=e\}}
\]
refinan las dos medidas espectrales diagonales y recuperan sus marginales. El
resultado no prueba inyectividad conjunta ni obliga a que un lector sea
compresión del otro. En el régimen no conmutativo, dos PVM qutrit mutuamente
insesgadas no poseen una PVM conjunta nítida; la realización simultánea exige
ruido compatible o un ambiente mayor. Esta diferencia de tipo preserva la
doble proyección HMT: el teorema matricial clasifica realizaciones de lectores
y no reescribe su antecedente fundacional.

\subsection*{Clausura física: dominio común, dos orientaciones y localización}

La cuarta clausura construye el dominio desde el que una ruta TPK puede recibir
simultáneamente realización espectral, hilbertiana, dimensional y de Cartan.
Sean \(\mathsf D_{\rm sp}\), \(\mathsf D_Q\), \(\mathsf D_D\) y
\(\mathsf D_C\) los dominios de referencia de esas cuatro representaciones, con
funtores de olvido fieles hacia la categoría
\(\mathsf P_{\rm TPK}^{\rm enr}\) de estados preparados y rutas registradas.
El producto fibrado
\[
 \mathsf P_{\rm CT}^{+}
 :=\mathsf D_{\rm sp}
 \times_{\mathsf P_{\rm TPK}^{\rm enr}}\mathsf D_Q
 \times_{\mathsf P_{\rm TPK}^{\rm enr}}\mathsf D_D
 \times_{\mathsf P_{\rm TPK}^{\rm enr}}\mathsf D_C
\]
impone que las cuatro componentes tengan exactamente el mismo antecedente.
Sus proyecciones canónicas son adaptadores functoriales y producen
\[
 \mathfrak R_{\rm MD}^{\rm disc,+}
 =\bigl(\mathfrak S_+,\mathcal Q_+,\mathfrak D_+,
        \rho_{C,+}^{\rm disc}\bigr),
\]
con valores, respectivamente, en rutas espectrales, espacios de Hilbert reales
con estructura compleja, el monoide dimensional positivo y el grupo de
Poincaré discreto. La construcción afirma el funtor sobre el dominio común; no
afirma que toda ruta TPK admita automáticamente los cuatro levantamientos.

Sobre una ruta \(\gamma=e_1\cdots e_r\), la variación total
\[
 V_+(\gamma)=\sum_jw(e_j)
\]
y el transporte neto orientado
\[
 V_{\rm or}(\gamma)=\sum_j\varepsilon(e_j)w(\underline e_j)
\]
son lectores hermanos. El lazo \(ee^\vee\) demuestra que el primero no
factoriza, en general, por el segundo: posee transporte orientado nulo y
variación positiva. Para obtener una realización invertible se restringe el
dominio a las aristas cuyas componentes espectral e hilbertiana son
invertibles, cuyo dato dimensional es una cocadena aditiva y cuyos cociclos de
Cartan son normalizados y antisimétricos. Sólo entonces la localización
\(\Pi_\vee(\mathsf P_{{\rm CT},{\rm rev}}^+)\) produce el funtor orientado
\(\mathfrak R_{\rm MD}^{\rm disc,or}\) y transporta la inversión componente a
componente. La lectura positiva no se elimina y la orientada no se obtiene
cambiando únicamente un signo.

Esta construcción física precede a las construcciones de Planck y del electrón
porque les proporciona dominio, representaciones y reglas de transporte, pero
no absorbe sus demostraciones de referencia ni la Ley General de Masas. La
cadena número--partícula--ruta--registro cronológico enriquecido--firma--
carácter--torres conserva su dirección causal; la escala absoluta de acción,
el electrón central, las familias de masa, CKM, PMNS y los sectores físicos
mantienen sus objetos, dominios y conclusiones propios. Una coincidencia metrológica se
contrasta sólo después de congelar la salida interna. Las extensiones que aún
requieren un entrelazador, un límite suave, una prueba de estabilidad dinámica
o un morfismo físico se enuncian con esos datos como hipótesis exactas. La
ausencia de alguno de esos mapas impide extender el teorema al codominio
correspondiente, sin alterar los teoremas genealógicos, operatoriales y
matriciales ya demostrados en sus dominios.

La lectura de la obra sigue, en consecuencia, el orden de dependencia y no el
de reconocimiento exterior:
\[
\begin{aligned}
 \mathrm{APP}&\longrightarrow\mathrm{TRIT}_{\rm orientado}
 \longrightarrow\mathrm{TPK}\\
 &\longrightarrow\text{sección enriquecida}
 \longrightarrow
 \begin{cases}
  \text{evaluación arquimediana},\\
  \text{proyección excepcional},\\
  \text{álgebras cilíndricas y clausura tracial},\\
  \text{rutas preparadas y realización física}.
 \end{cases}
\end{aligned}
\]
Cada salida conserva su dominio, sus hipótesis y sus obstrucciones precisas. Lo común no es
una igualdad forzada entre codominios, sino la procedencia desde el mismo
estado enriquecido. Éste es el sentido preciso del cierre holográfico: ninguna
publicación terminal agota la genealogía que la produce y ninguna clausura
posterior puede sustituir la doble dirección --arquimediana y excepcional--
que conserva la unidad con memoria.
